\documentclass[10pt,a4paper]{amsart}
\usepackage[margin=19mm]{geometry}
\usepackage{amsmath,amssymb,mathtools}
\usepackage[scr=rsfs,cal=euler]{mathalfa}
\usepackage{xfrac}
\usepackage[unicode,hidelinks,bookmarks=false]{hyperref}
\newcommand{\ie}{i.e.\ }
\newcommand{\cf}{cf.\ }
\newcommand{\resp}{respectively\ }
\newcommand{\Subsection}{\subsection}
\providecommand{\nobreakdash}{\nobreak\hbox{-}}
\setlength{\emergencystretch}{2em}
\multlinegap0pt
\usepackage{bm}
\usepackage{bbm}
\usepackage{dsfont}
\usepackage{xspace}
\usepackage{cancel}
\usepackage{mathtools}
\usepackage{amsthm}
\makeatletter
\@ifundefined{theorem}{\newtheorem{theorem}{Theorem}}{}
\@ifundefined{corollary}{\newtheorem{corollary}[theorem]{Corollary}}{}
\@ifundefined{lemma}{\newtheorem{lemma}[theorem]{Lemma}}{}
\makeatother
\newcommand{\ZR}{\mathbb{R}}
\newcommand{\cl}{\mathcal L}
\newcommand{\cs}{\mathcal S}
\newcommand{\de}{\delta}
\title[A Kakeya maximal estimate for regulus strips]{A Kakeya maximal estimate for regulus strips}
\author{Shukun Wu}
\date{Source: arXiv:2411.04438v2 (CC BY 4.0)}
\begin{document}
\maketitle

\noindent\emph{Source and licence.} A Kakeya maximal estimate for regulus strips. Version arXiv:2411.04438v2 (CC BY 4.0), distributed under the Creative Commons Attribution 4.0 International licence, as recorded in the arXiv licence metadata for this exact version and shown on the abstract page https://arxiv.org/abs/2411.04438v2. Full rights notice: licenses/arxiv-2411.04438-CC-BY-4.0.txt.\par\smallskip

\noindent\textit{Editorial scope.} Counted unit: the Corollary whose source label is \texttt{regulus-strip-repara} in the article (source0.tex lines 582--596, source label \texttt{regulus-strip-repara}), delivered together with the article's own complete original proof of it (source0.tex lines 597--661, one \texttt{proof} environment of 65 lines). No mathematics is written, repaired or re-derived: statement and proof are the article's own text line for line. Required context retained uncounted: regulus-first-para (lines 369--372); family-regulus-strip (lines 393--396); sl2-is-linear (lines 535--538). Every definition, display and label of those lines is retained unchanged. Omitted from this package and not redistributed: 1--368, 373--392, 397--534, 539--581, 662--1320. Nothing inside a retained excerpt is omitted or altered: every retained body is delivered exactly as the article's lines are written, so no omission receipt of any kind accompanies this package. Citations: the retained excerpts carry no citation command at all, so no bibliography entry of the article is transcribed and no reference is redistributed. Reference adaptation: none was needed and none was made; every reference inside the delivered unit resolves inside it, so no unresolved reference or citation is produced. Printed slips kept literal: no printed word, symbol, hypothesis or number of the counted statement or of its proof was altered. Layout and class: the article's own class is \texttt{amsart} with packages including tikz, caption, xcolor, amssymb, latexsym, amsmath, extarrows, graphicx, mathrsfs, comment, hyperref, url, pict2e, amstext; the wrapper is the lane's pinned \texttt{amsart} editorial wrapper at 10pt on A4, which restates the theorem-like environments the retained text needs and adds the article's own macro definitions that the retained text needs (\texttt{\textbackslash ZR}, \texttt{\textbackslash cl}, \texttt{\textbackslash cs}, \texttt{\textbackslash de}), with extra wrapper packages (bm, bbm, dsfont, xspace, cancel, mathtools). The delivered numbering is the unit's own. Assets: no retained span contains a figure environment, a graphics-inclusion command, a PDF-page inclusion, a table or a TikZ picture; no image file is copied into this package, the package declares no asset dependency and no image file is redistributed with it. Licence: version arXiv:2411.04438v2 of this article is distributed under the Creative Commons Attribution 4.0 International licence, as recorded in the arXiv licence metadata for this exact version and shown on the abstract page https://arxiv.org/abs/2411.04438v2. A full rights notice is delivered at licenses/arxiv-2411.04438-CC-BY-4.0.txt. Characters: every retained excerpt is pure ASCII, so the delivered engine, XeLaTeX with its default Latin Modern text font, reports no missing character.
\begin{equation}
\label{regulus-first-para}
    R(\ell):=\bigcup_{p\in\ell}\ell(p).
\end{equation}

\begin{equation}
\label{family-regulus-strip}
    \bar\cs:=\{S(\ell_{(a,b,c,d)}):a,d\sim1,ad-bc=1\}.
\end{equation}

\begin{lemma}[$\rm{SL}_2$ is linear]
\label{sl2-is-linear}
Lines in $\cl_{\text{SL}_2}$ satisfying $ad\not=0$ can be identified as lines in $\cl$.
\end{lemma}

\begin{corollary}
\label{regulus-strip-repara}
Let $\bar\cs$ be the family of regulus strips defined in \eqref{family-regulus-strip}.
After choosing an appropriate parameterization, each $S(\ell)\in \bar\cs$ can be expressed as
\begin{equation}
    S(\ell)=\{(a+ct, b+at, t)+u(1,-t,0), \hspace{.3cm}t\in[0,1], \,\,\,|u|\lesssim\de^{1/2}\}+B_\de^3,
\end{equation}
where $\ell=\ell_{(a,b,c)}\in\cl$.
In particular, if we let 
\begin{equation}
\label{regulus}
    R(\ell_{(a,b,c)}):=\{(a+ct, b+at, t)+u(1,-t,0), \hspace{.3cm}t\in[0,1], \,\,\,u\in\ZR\}
\end{equation}
be a regulus, then $S(\ell)=N_\de(R(\ell))\cap N_{\de^{1/2}}(\ell)$.
\end{corollary}

\begin{proof}
Suppose $\ell$ is parameterized as $\ell_{(A,B,C,D)}$ with $AD-BC=1$.
By \eqref{regulus-first-para}, we can write the regulus $R(\ell)$ explicitly as
\begin{align}
\label{regulus-para-explicit}
    R(\ell)&=\left\{(A+Cs,B+Ds,s)+r(A+Cs,B+Ds,0):    s,r\in\mathbb R\right\}\\
    &=\left\{
    (\lambda(A+Cs),\lambda(B+Ds),s):
    s,\lambda\in\mathbb R
    \right\}.
\end{align}


As in Lemma \ref{sl2-is-linear}, consider the coordinate swap $\Phi(X,Y,Z)=(X,Z,Y)$.
The new height variable is $t=Y$. 
Let
\begin{equation}
    a=\frac1D,\qquad b=-\frac BD,\qquad c=\frac CD.
\end{equation}
Note that a point of the regulus $R(\ell)$ is $(\lambda(A+Cs),\lambda(B+Ds),s)$.
After the coordinate swap, its new height is $t=\lambda(B+Ds)$.
Solving for $s$, we get
\begin{equation}
    s=\frac{t-\lambda B}{\lambda D}
      =
      -\frac BD+\frac{t}{\lambda D}
      =
      b+\frac a\lambda t.
\end{equation}
As a result, the first coordinate of $R(\ell)$ becomes
\begin{equation}
        \lambda(A+Cs)
    =
    \lambda A+\lambda C\frac{t-\lambda B}{\lambda D}
    =
    \lambda\left(A-\frac{CB}{D}\right)+\frac CD t=a\lambda+ct,
\end{equation}
since $AD-BC=1$.
Therefore, after the coordinate swap, the regulus $\Phi(R(\ell))$ can be parameterized as
\begin{equation}
    \left\{
    \left(a\lambda+ct,\ b+\frac a\lambda t,\ t\right):
    t,\lambda\in\mathbb R
    \right\}.
\end{equation}

\smallskip 

Put $ u=a(\lambda-1)$, so $\lambda=1+\frac ua$.
Consequently,
\begin{equation}
    \left(a\lambda+ct,\ b+\frac a\lambda t,\ t\right)
    =
    (a+ct,b+at,t)+u(1,-t,0)
    +
    \left(0,\frac{u^2}{a+u}t,0\right)
\end{equation}
Note that the assumptions $|u|\lesssim \delta^{1/2}$ and $a\sim 1$ imply that $\frac{u^2}{a+u}t=O(\delta)$
for $t\in[0,1]$. 
This completes the proof of the Corollary.
\qedhere



\end{proof}

\end{document}
